\section{Early Stopping}
\paragraph{Purpose}
With a model large enough to overfit, the training error keeps decreasing while the
validation error eventually starts to rise again. \tB{Stop training when the validation
error stops improving} and return the parameters of the best epoch.

\paragraph{Theory}
\begin{enumerate}
    \item evaluate the validation error every few steps;
    \item keep a copy of the parameters each time it improves;
    \item stop after $\text{patience}$ evaluations without improvement and return the
        saved parameters.
\end{enumerate}
The number of training steps is treated as a hyper-parameter chosen on the validation
set. For a quadratic objective with gradient descent, it can be shown to be equivalent to
$L^{2}$ regularization, the number of steps playing the role of $\frac{1}{\alpha}$.

\paragraph{Strengths}
Very cheap, requires no change to the model or to the objective.

\paragraph{Weaknesses}
Requires held-out data. Once the number of steps is known, the model can be retrained on
the training and validation sets together for the same number of steps.
